\section{Tutorial 6: Optimization Behavior}

\subsection{Objectives}

This tutorial develops intuition for how optimization behaves in practice:

\begin{itemize}
    \item Understand the behavior of gradient descent vs SGD
    \item Analyze the effect of learning rate and batch size
    \item Study conditioning and convergence speed
    \item Connect optimization dynamics to generalization
\end{itemize}

\medskip

\noindent
\textbf{Focus.}  
Go beyond formulas and understand how optimization behaves in real training.

\subsection{Exercise 1: Gradient Descent on a Quadratic}

Consider:
\[
f(w) = \frac{1}{2} w^\top A w,
\]
where $A$ is symmetric positive definite.

\medskip

\noindent
\textbf{Tasks.}
\begin{enumerate}
    \item Write the gradient descent update.
    \item Describe convergence behavior in terms of eigenvalues of $A$.
\end{enumerate}

\medskip

\noindent
\textbf{Solution.}

\begin{enumerate}
\item Update:
\[
w_{t+1} = w_t - \eta A w_t.
\]

\item Let eigenvalues of $A$ be $\lambda_1 \leq \cdots \leq \lambda_d$.

\begin{itemize}
    \item Convergence rate depends on condition number:
    \[
    \kappa = \frac{\lambda_{\max}}{\lambda_{\min}}.
    \]
    \item Large $\kappa$ $\Rightarrow$ slow convergence
\end{itemize}
\end{enumerate}

\medskip

\noindent
\textbf{Insight.}  
Ill-conditioning leads to slow and anisotropic convergence.

\subsection{Exercise 2: Effect of Learning Rate}

Consider gradient descent with step size $\eta$.

\medskip

\noindent
\textbf{Tasks.}
\begin{enumerate}
    \item What happens if $\eta$ is too small?
    \item What happens if $\eta$ is too large?
\end{enumerate}

\medskip

\noindent
\textbf{Solution.}

\begin{enumerate}
\item Small $\eta$:
\begin{itemize}
    \item Slow convergence
\end{itemize}

\item Large $\eta$:
\begin{itemize}
    \item Oscillations or divergence
\end{itemize}
\end{enumerate}

\medskip

\noindent
\textbf{Key takeaway.}  
Learning rate critically controls stability and speed.

\subsection{Exercise 3: SGD vs Full Gradient}

Consider minimizing:
\[
L(\theta) = \frac{1}{n} \sum_{i=1}^n \ell_i(\theta).
\]

\medskip

\noindent
\textbf{Tasks.}
\begin{enumerate}
    \item Compare full gradient descent and SGD updates.
    \item Why does SGD introduce noise?
\end{enumerate}

\medskip

\noindent
\textbf{Solution.}

\begin{enumerate}
\item 
\begin{itemize}
    \item GD uses full gradient
    \item SGD uses a single sample or mini-batch
\end{itemize}

\item Noise arises because:
\[
\nabla \ell_i(\theta) \neq \nabla L(\theta)
\]
for a single sample.
\end{enumerate}

\medskip

\noindent
\textbf{Insight.}  
SGD trades accuracy of gradients for efficiency.

\subsection{Exercise 4: Role of Batch Size}

\textbf{Tasks.}
\begin{enumerate}
    \item What happens as batch size increases?
    \item How does batch size affect generalization?
\end{enumerate}

\medskip

\noindent
\textbf{Solution.}

\begin{enumerate}
\item Larger batch:
\begin{itemize}
    \item Lower variance gradient
    \item More stable updates
\end{itemize}

\item Smaller batch:
\begin{itemize}
    \item More noise
    \item Often better generalization
\end{itemize}
\end{enumerate}

\medskip

\noindent
\textbf{Insight.}  
Noise in SGD can act as implicit regularization.

\subsection{Exercise 5: Momentum Interpretation}

Consider momentum updates:
\[
v_{t+1} = \beta v_t + \nabla L(\theta_t).
\]

\medskip

\noindent
\textbf{Tasks.}
\begin{enumerate}
    \item Explain the role of $\beta$.
    \item Why does momentum help in ill-conditioned problems?
\end{enumerate}

\medskip

\noindent
\textbf{Solution.}

\begin{enumerate}
\item $\beta$ controls how much past gradients influence updates.

\item Momentum:
\begin{itemize}
    \item Accelerates movement in consistent directions
    \item Reduces oscillations in narrow valleys
\end{itemize}
\end{enumerate}

\medskip

\noindent
\textbf{Insight.}  
Momentum smooths optimization trajectories.

\subsection{Exercise 6: Implicit Regularization of SGD}

\textbf{Question.}  
Why might SGD generalize better than deterministic gradient descent?

\medskip

\noindent
\textbf{Discussion points.}

\begin{itemize}
    \item Noise helps escape sharp minima
    \item Prefers flatter regions of the loss landscape
    \item Acts as implicit regularization
\end{itemize}

\medskip

\noindent
\textbf{Insight.}  
Optimization dynamics influence the solutions found.

\subsection{Exercise 7: Early Stopping as Regularization}

\textbf{Tasks.}
\begin{enumerate}
    \item Why can training too long lead to overfitting?
    \item How does early stopping help?
\end{enumerate}

\medskip

\noindent
\textbf{Solution.}

\begin{enumerate}
\item Long training:
\begin{itemize}
    \item Model fits noise
\end{itemize}

\item Early stopping:
\begin{itemize}
    \item Limits effective complexity
\end{itemize}
\end{enumerate}

\medskip

\noindent
\textbf{Insight.}  
Training time controls model complexity.

\subsection{Exercise 8 (Discussion): Flat vs Sharp Minima}

\textbf{Question.}  
Why are flat minima often associated with better generalization?

\medskip

\noindent
\textbf{Discussion points.}

\begin{itemize}
    \item Robustness to perturbations
    \item Stability under noise
    \item Relation to SGD dynamics
\end{itemize}

\medskip

\noindent
\textbf{Insight.}  
Geometry of the loss landscape affects generalization.

\subsection{Summary}

\begin{itemize}
    \item Optimization behavior depends on learning rate, batch size, and conditioning
    \item SGD introduces noise that affects both convergence and generalization
    \item Momentum and adaptive methods improve optimization efficiency
    \item Optimization dynamics act as implicit regularization
\end{itemize}

\medskip

\noindent
\textbf{Preparation for next lecture.}  
We now study numerical and statistical stability, examining how scaling and noise affect learning behavior.